\section{Задание}
Доработать ИС из ЛР1 следующим образом:


1. Добавить в систему возможность массового добавления объектов при помощи импорта файла. Формат для импорта необходимо согласовать с преподавателем. Импортируемый файл должен загружаться на сервер через интерфейс разработанного веб-приложения.


- При реализации логики импорта объектов необходимо реализовать транзакцию таким образом, чтобы в случае возникновения ошибки при импорте не был создан ни один объект.


- При импорте должна быть реализована проверка пользовательского ввода в соответствии с ограничениями предметной области из ЛР1.


- При наличии вложенных объектов в основном объекте из ЛР1 необходимо задавать значения полей вложенных объектов в той же записи, что и основной объект.


2. Необходимо добавить в систему интерфейс для отображения истории импорта (обычный пользователь видит только операции импорта, запущенные им, администратор — все операции).


- В истории должны отображаться ID операции, статус её завершения, пользователь, который её запустил, число добавленных объектов в операции (только для успешно завершённых).


3. Согласовать с преподавателем и добавить в модель из первой лабораторной новые ограничения уникальности, проверяемые на программном уровне (эти новые ограничения должны быть реализованы в рамках бизнес-логики приложения и не должны быть отображены/реализованы в БД).


4. Реализовать сценарий с использованием Apache JMeter, имитирующий одновременную работу нескольких пользователей с ИС, и проверить корректность изоляции транзакций, используемых в ЛР. По итогам исследования поведения системы при её одновременном использовании несколькими пользователями изменить уровень изоляции транзакций там, где это требуется. Обосновать изменения.


- Реализованный сценарий должен покрывать создание, редактирование, удаление и импорт объектов.


- Реализованный сценарий должен проверять корректность поведения системы при попытке нескольких пользователей обновить и/или удалить один и тот же объект (например, двух администраторов).


- Реализованный сценарий должен проверять корректность соблюдения системой ограничений уникальности предметной области при одновременной попытке нескольких пользователей создать объект с одним и тем же уникальным значением.


\subsection{Требования к отчёту}
\begin{enumerate}\setlength{\itemsep}{0pt}
\item Текст задания.
\item UML-диаграммы классов и пакетов разработанного приложения.
\item Исходный код системы или ссылка на репозиторий с исходным кодом.
\item Выводы по работе.
\end{enumerate}

\subsection{Вопросы к защите}
\begin{enumerate}\setlength{\itemsep}{0pt}
\item Понятие бизнес-логики в программных системах. Уровень бизнес-логики в многоуровневой архитектуре программных систем.
\item Jakarta Enterprise Beans (EJB). Виды бинов и их назначение.
\item EJB Session beans. Жизненный цикл.
\item Понятие транзакции. Транзакции в БД. ACID.
\item Виды конфликтов при многопользовательской работе с данными. Уровни изоляции транзакций.
\item Особенности реализации транзакций на уровне бизнес-логики, отличия от транзакций на уровне БД.
\item Java Transaction API. Основные принципы и программные интерфейсы.
\item Реализация управления транзакциями в Jakarta EE. Декларативное и программное управление транзакциями.
\item Реализация управления транзакциями в Spring. Декларативное и программное управление транзакциями в Spring. Аннотация \texttt{@Transactional}.
\end{enumerate}

\subsection{Принятые решения}
С преподавателем согласованы формат JSON и уникальность сочетания названия фильма и режиссёра, реализуемая только в бизнес-логике приложения. Система развивает приложение первой лабораторной работы и сохраняет операции CRUD и специальные операции.